%!TEX root = ../template.tex
%%%%%%%%%%%%%%%%%%%%%%%%%%%%%%%%%%%%%%%%%%%%%%%%%%%%%%%%%%%%%%%%%%%%
%% chapter7.tex
%% NOVA thesis document file
%%
%% Chapter with lots of dummy text
%%%%%%%%%%%%%%%%%%%%%%%%%%%%%%%%%%%%%%%%%%%%%%%%%%%%%%%%%%%%%%%%%%%%

\typeout{NT FILE chapter7.tex}%

\chapter{Comparative Evaluation}
\label{cha:Comparative_Evaluation}

In order to create a fair comparison for both verification approaches the same use case will be used and analysed thoroughly for the
extraction approach and translation approach. Before going into the results it is fundamental to talk about the impact of the subset
in choosing the examples that can or not be used from the various verified examples from cameleer. The use case should, as always,
respect the limitations implied from the subset chosen.


\section{Verification Guarantees and Trust Boundaries}
\label{sec:Verification_Guarantees_and_Trust_Boundaries}

Before presenting the comparative results, it is necessary to characterize the trust boundaries of each approach and identify where verification 
guarantees can break down. Both approaches rely on an unverified component that constitutes the critical point of failure in the formal verification.

In the extraction approach, the code is verified within \whythree and the proof obligations are discharged before extraction is invoked as seen in figure \ref{fig:Cameleer_pipeline}. 
However, the extraction driver itself is not formally verified, meaning that if it fails to faithfully reproduce the \whyml semantics in the generated \cml code, 
the verification guarantees established in \whythree are silently lost. In this scenario, even if the resulting \cml code compiles successfully 
through the verified compiler, the compiler's correctness theorem offers no assurance that the verified properties are preserved, meaning the guarantee 
only extends from the \cml source downwards, not from the \whyml source through the driver.

For the translation approach, the failure point arises earlier in the pipeline shown in figure \ref{fig:GoalPipeline}. If the translation tool incorrectly translates \cml with 
\gospel specification into \whyml, the resulting \whyml program may not faithfully represent the original \cml code, causing the verification step in \whythree 
to verify an inaccurate model. In this case, even if the proofs discharge successfully, they would be proving properties of the translated model rather than 
the original program. On the other hand, if the translation is faithful and the proofs discharge, the verified \cml compiler provides a strong guarantee that the 
verified properties are preserved all the way to machine code.

The fundamental distinction between the two approaches therefore lies not in the verification mechanism itself, which is \whythree in both cases, but in the position 
of the unverified component and the consequences of its failure.

\section{Comparative Case Studies}

The case studies explored in these chapter were selected accordingly with the subset of the \whyml. For each example, the same 
program is subjected to both approaches independently, allowing a direct comparison of the verification effort required, the 
expressiveness of the resulting specifications, and the degree to which each method preserves the intent of the original code.

\subsection{Red-Black Tree}

The Red-Black tree is a classical example of a self-balancing binary search tree in which each node is assigned a colour, either 
Red or Black, subject to two structural invariants: no two adjacent nodes may both be Red, and every path from a node to a 
leaf must contain the same number of Black nodes~\cite{Guibas1978}. These invariants together guarantee that the height of the tree remains bounded 
by \textit{O}(log \textit{n}), ensuring efficient worst-case performance for search, insertion, and deletion. Furthermore, if 
the tree as a whole satisfies these conditions, every subtree is also guaranteed to satisfy them, making the invariant compositional by nature.

The Red-Black tree data structure has a data type of tree with it being a leaf or a node that has the values of the color,
the key, the value and both left and right trees.

\begin{whylang}
    type key = int
    type value = int

    type color = Red | Black

    type tree =
        | Leaf
        | Node color tree key value tree
\end{whylang}

Before going into the functions there is the need to create some predicates and lemmas to verify those functions. 
The predicate \inlinecode{memt}
checks whether a key value pair \inlinecode{(k, v)} occurs anywhere in the tree \inlinecode{t}, if it is a \inlinecode{Leaf} then it returns false, 
if \inlinecode{t} is a \inlinecode{Node} then it verifies if the key and value presented to \inlinecode{memt} are the same as the key and value as
\inlinecode{t} or it checks if the pair appears in both other subtrees.

\begin{whylang}
    predicate memt (t : tree) (k : key) (v : value) =
        match t with
        | Leaf -> false
        | Node _ l k' v' r -> (k = k' /\ v = v') \/ memt l k v \/ memt r k v
        end
\end{whylang}

The lemma \inlinecode{memt\_color} states that membership is independent of \inlinecode{Node} color.

\begin{whylang}
    lemma memt_color :
        forall l r : tree, k k' : key, v v' : value, c c' : color.
        memt (Node c l k v r) k' v' -> memt (Node c' l k v r) k' v'
\end{whylang}

Predicate \inlinecode{lt\_tree} and \inlinecode{gt\_tree} each state that every key in tree \inlinecode{t} is strictly lesser than \inlinecode{x} and 
that every key in tree \inlinecode{t} is strictly greater than \inlinecode{x} respectively, this is to express the left and right tree sub-conditions 
respectively.

\begin{whylang}
    predicate lt_tree (x : key) (t : tree) =
        forall k : key. forall v : value. memt t k v -> k < x

    predicate gt_tree (x : key) (t : tree) =
        forall k : key. forall v : value. memt t k v -> x < k
\end{whylang}

The trivial lemmas \inlinecode{lt\_leaf} and \inlinecode{gt\_leaf} are true base cases for an empty tree that satisfies both.

\begin{whylang}
    lemma lt_leaf: forall x: key. lt_tree x Leaf

    lemma gt_leaf: forall x: key. gt_tree x Leaf
\end{whylang}

The construction lemmas \inlinecode{lt\_tree\_node} and \inlinecode{gt\_tree\_node} state that if both subtrees and the node key
satisfies the ordering condition then the complete node also satisfies that condition.

\begin{whylang}
    lemma lt_tree_node:
        forall x y: key, v: value, l r: tree, c: color.
        lt_tree x l -> lt_tree x r -> y < x -> lt_tree x (Node c l y v r)

    lemma gt_tree_node:
        forall x y: key, v: value, l r: tree, c: color.
        gt_tree x l -> gt_tree x r -> x < y -> gt_tree x (Node c l y v r)
\end{whylang}

The extraction lemma \inlinecode{lt\_node\_lt} says that if the whole node satisfies the \inlinecode{lt\_tree x} then the key of
node needs to be lesser than \inlinecode{x} and similarly for lemma \inlinecode{gt\_node\_gt} but if it satifies \inlinecode{gt\_tree x}
then the key needs to be greater.

\begin{whylang}
    lemma lt_node_lt:
        forall x y: key, v: value, l r: tree, c: color.
        lt_tree x (Node c l y v r) -> y < x

    lemma gt_node_gt:
        forall x y: key, v: value, l r: tree, c: color.
        gt_tree x (Node c l y v r) -> x < y
\end{whylang}

The propagation lemmas \inlinecode{lt\_left}, \inlinecode{lt\_right}, \inlinecode{gt\_right} and \inlinecode{gt\_right} says that
if a node satisfies either \inlinecode{lt\_tree x} or \inlinecode{gt\_tree x} then both left and right subtrees need to satisfy
the same conditions.

\begin{whylang}
    lemma lt_left:
        forall x y: key, v: value, l r: tree, c : color.
        lt_tree x (Node c l y v r) -> lt_tree x l

    lemma lt_right:
        forall x y: key, v: value,  l r: tree, c: color.
        lt_tree x (Node c l y v r) -> lt_tree x r

    lemma gt_left:
        forall x y: key, v: value, l r: tree, c: color.
        gt_tree x (Node c l y v r) -> gt_tree x l

    lemma gt_right:
        forall x y: key, v: value, l r: tree, c: color.
        gt_tree x (Node c l y v r) -> gt_tree x r
\end{whylang}

The exclusion lemma \inlinecode{lt\_tree\_not\_in} state that if all the keys in a tree is less that \inlinecode{x}, then \inlinecode{x} 
cannot be a key for the tree and for lemma \inlinecode{gt\_tree\_not\_in} is the opposite, if \inlinecode{x} is greater than it cannot be a key.

\begin{whylang}
    lemma lt_tree_not_in:
        forall x: key, t: tree. lt_tree x t -> forall v: value. not (memt t x v)

    lemma gt_tree_not_in :
        forall x: key, t: tree. gt_tree x t -> forall v: value. not (memt t x v)
\end{whylang}

The transitivity lemma \inlinecode{lt\_tree\_trans} says that if all keys in a tree is less than \inlinecode{x} and \inlinecode{x < y}, then all
keys are also less than \inlinecode{y}, the same goes for lemma \inlinecode{gt\_tree\_trans} but if all the keys are greater than \inlinecode{x}
and \inlinecode{y < x}, then all keys are greater than \inlinecode{y}.

\begin{whylang}
    lemma lt_tree_trans:
        forall x y: key. x < y -> forall t: tree. lt_tree x t -> lt_tree y t

    lemma gt_tree_trans :
        forall x y: key. y < x -> forall t: tree. gt_tree x t -> gt_tree y t
\end{whylang}

Predicate \inlinecode{bst} represents the invariant of the search tree, where every \inlinecode{Leaf} is a \inlinecode{bst} and
a \inlinecode{Node} is \inlinecode{bst} if both subtrees are \inlinecode{bst}, if all the keys from the left subtree are less than the
node tree and all the keys from the right subtree are greater.

\begin{whylang}
    predicate bst (t : tree) =
        match t with
        | Leaf -> true
        | Node _ l k _ r -> bst l /\ bst r /\ lt_tree k l /\ gt_tree k r
        end
\end{whylang}

The lemma \inlinecode{bst\_Leaf} is a simple representation of a \inlinecode{Leaf} that is a \inlinecode{bst}.

\begin{whylang}
    lemma bst_Leaf : bst Leaf
\end{whylang}

The following lemmas represent that if a node is a \inlinecode{bst} than for each subtree left or right it also must be a \inlinecode{bst}.

\begin{whylang}
    lemma bst_left:
        forall k: key, v: value, l r: tree, c: color. bst (Node c l k v r) -> bst l

    lemma bst_right:
        forall k: key, v: value, l r: tree, c: color. bst (Node c l k v r) -> bst r
\end{whylang}

The lemma \inlinecode{bst\_color} says that \inlinecode{bst} condition is independent of color.

\begin{whylang}
    lemma bst_color:
        forall c c': color, k: key, v: value, l r: tree.
        bst (Node c l k v r) -> bst (Node c' l k v r)
\end{whylang}

The lemmas \inlinecode{rotate\_left} and \inlinecode{rotate\_right} represent the change of color and rotating the subtrees
while still satisfying the condition of \inlinecode{bst}.

\begin{whylang}
    lemma rotate_left:
        forall kx ky: key, vx vy: value, a b c: tree, c1 c2 c3 c4: color.
        bst (Node c1 a kx vx (Node c2 b ky vy c)) ->
        bst (Node c3 (Node c4 a kx vx b) ky vy c)

    lemma rotate_right:
        forall kx ky: key, vx vy: value, a b c: tree, c1 c2 c3 c4: color.
        bst (Node c3 (Node c4 a kx vx b) ky vy c) ->
        bst (Node c1 a kx vx (Node c2 b ky vy c))
\end{whylang}

The predicate to represent a tree that is not red is the following.

\begin{whylang}
    predicate is_not_red (t : tree) =
        match t with
        | Node Red _ _ _ _ -> false
        | Leaf | Node Black _ _ _ _ -> true
        end
\end{whylang}

The predicate that represents the red black tree is \inlinecode{rbtree} that returns zero if tree \inlinecode{t} is 
a \inlinecode{Leaf}. If the tree is a red node than it does the \inlinecode{rbtree} for both the left and right
subtrees without changing the value \inlinecode{n} received, while also making sure those subtrees are not red. Lastly,
if the node is a black node than it does \inlinecode{rbtree} while subtracting one from \inlinecode{n}.


\begin{whylang}
    predicate rbtree (n : int) (t : tree) =
        match t with
        | Leaf ->
            n = 0
        | Node Red   l _ _ r ->
            rbtree n l /\ rbtree n r /\ is_not_red l /\ is_not_red r
        | Node Black l _ _ r ->
            rbtree (n-1) l /\ rbtree (n-1) r
        end
\end{whylang}

The lemma \inlinecode{rbtree\_Leaf} is the representation of a \inlinecode{Leaf} that satisfies the predicate \inlinecode{rbtree}.

\begin{whylang}
    lemma rbtree_Leaf: rbtree 0 Leaf
\end{whylang}

Lemma \inlinecode{rbtree\_Node1} tells us that a red node that has both of the subtrees as leafs will be a \inlinecode{rbtree} of value
zero.

\begin{whylang}
    lemma rbtree_Node1:
        forall k:key, v:value. rbtree 0 (Node Red Leaf k v Leaf)
\end{whylang}

The lemmas \inlinecode{rbtree\_left} and \inlinecode{rbtree\_right} say that if a tree \inlinecode{t} is a \inlinecode{rbtree}
of value \inlinecode{n} than there exists a node in the subtree that also satisfies that \inlinecode{rbtree} condition.

\begin{whylang}
    lemma rbtree_left:
        forall x: key, v: value, l r: tree, c: color.
        (exists n: int. rbtree n (Node c l x v r)) -> exists n: int. rbtree n l

    lemma rbtree_right:
        forall x: key, v: value, l r: tree, c: color.
        (exists n: int. rbtree n (Node c l x v r)) -> exists n: int. rbtree n r
\end{whylang}

The predicate \inlinecode{almost\_rbtree} is an intermediate representation of a red-black tree that does not verify if the subtrees of
a red node are not red, this is used for when there is the need to satisfy the condition of red-black tree while it is doing some changes
to the tree.

\begin{whylang}
    predicate almost_rbtree (n : int) (t : tree) =
        match t with
        | Leaf ->
            n = 0
        | Node Red   l _ _ r ->
            rbtree n l /\ rbtree n r
        | Node Black l _ _ r ->
            rbtree (n-1) l /\ rbtree (n-1) r
        end
\end{whylang}

Lemma \inlinecode{rbtree\_almost\_rbtree} says that if a tree is \inlinecode{rbtree} than it is also \inlinecode{almost\_rbtree}.

\begin{whylang}
    lemma rbtree_almost_rbtree:
        forall n: int, t: tree. rbtree n t -> almost_rbtree n t
\end{whylang}

The following lemma \inlinecode{rbtree\_almost\_rbtree\_ex} says that for every tree there exists a value for the \inlinecode{rbtree} that
is the same as the \inlinecode{almost\_rbtree}.

\begin{whylang}
    lemma rbtree_almost_rbtree_ex:
        forall s: tree.
        (exists n: int. rbtree n s) -> exists n: int. almost_rbtree n s
\end{whylang}

The final lemma \inlinecode{almost\_rbtree\_rbtree\_black} says that for all the black nodes that are \inlinecode{almost\_rbtree} of
\inlinecode{n} are also \inlinecode{rbtree} of \inlinecode{n}, which means that from the black trees standpoint there was no difference from
\inlinecode{rbtree} and \inlinecode{almost\_rbtree}.

\begin{whylang}
    lemma almost_rbtree_rbtree_black:
        forall x: key, v: value, l r: tree, n: int.
        almost_rbtree n (Node Black l x v r) -> rbtree n (Node Black l x v r)
\end{whylang}

After creating all the necessary gospel annotations we can correctly define the necessary invariants, pre-conditions and post-conditions for the
following functions supported in the red-black trees data structure.

The first function is the \inlinecode{find} that receives a tree \inlinecode{t} and a key \inlinecode{k} in which it traverses the tree recursively
until it finds the tree with the same key or until it reaches all the leafs and raises the exception \inlinecode{Not\_found}. This function requires
the \inlinecode{t} to be a \inlinecode{bst} and ensures the value returned is in line with the condition \inlinecode{memt}. It also needs to be specified
that \inlinecode{t} variates with each passing and that raises \inlinecode{Not\_found} if the key is not equal to a key of that tree.

\begin{whylang}
    let rec find (t : tree) (k : key) : value
        requires { bst t }
        variant { t }
        ensures { memt t k result }
        raises { Not_found -> forall v : value. not (memt t k v) }
    = match t with
        | Leaf -> raise Not_found
        | Node _ l k' v r ->
            if k = k' then v
            else if k < k' then find l k
            else (* k > k' *) find r k
        end
\end{whylang}

The functions need to balance each of the subtrees are \inlinecode{lbalance} for the left subtree and \inlinecode{rbalance}. The way they balance
the tree is to receive both subtrees, the value and the key, comparing what the subtree in which it was called this balancing function and rotating
the nodes they are not a leafs. 

The functions require the key \inlinecode{k} to be lesser than the other keys in left subtree and greater in the
right subtree, while the left and right subtrees satisfy the condition of \inlinecode{bst}. It also needs to ensure that the resulting tree is
satisfying the condition of \inlinecode{bst}, every value \inlinecode{n} of an \inlinecode{almost\_rbtree} of the respecting subtree, left for \inlinecode{lbalance}
and right for \inlinecode{rbalance}, implies that the opposing tree satisfies \inlinecode{rbtree} for the \inlinecode{n} and the resulting tree satifies
\inlinecode{bst} for that \inlinecode{n+1}. To ensure that the balancing was done correctly is to say that for every key \inlinecode{k'} and value \inlinecode{v\'}
that satisfies \inlinecode{memt} for the resulting tree than it is equivalent to satisfy \inlinecode{memt} for the left and right subtrees if the key \inlinecode{k\'}
is not the same as the key \inlinecode{k}, otherwise the values must be the same.

\begin{whylang}
    let lbalance (l : tree) (k : key) (v : value) (r : tree)
        requires { lt_tree k l /\ gt_tree k r /\ bst l /\ bst r }
        ensures { bst result /\
        (forall n : int. almost_rbtree n l -> rbtree n r -> rbtree (n+1) result)
        /\
        forall k':key, v':value.
            memt result k' v' <->
            if k' = k then v' = v else (memt l k' v' \/ memt r k' v') }
    = match l with
        | Node Red (Node Red a kx vx b) ky vy c
        | Node Red a kx vx (Node Red b ky vy c) ->
            Node Red (Node Black a kx vx b) ky vy (Node Black c k v r)
        | _ ->
            Node Black l k v r
        end

    let rbalance (l : tree) (k : key) (v : value) (r : tree)
        requires { lt_tree k l /\ gt_tree k r /\ bst l /\ bst r }
        ensures { bst result /\
        (forall n : int. almost_rbtree n r -> rbtree n l -> rbtree (n+1) result)
        /\
        forall k':key, v':value.
            memt result k' v' <->
            if k' = k then v' = v else (memt l k' v' \/ memt r k' v') }
    = match r with
        | Node Red (Node Red b ky vy c) kz vz d
        | Node Red b ky vy (Node Red c kz vz d) ->
            Node Red (Node Black l k v b) ky vy (Node Black c kz vz d)
        | _ ->
            Node Black l k v r
        end
\end{whylang}

The recursive function insert receives the \inlinecode{tree} and the key \inlinecode{k} and value \inlinecode{v} of the node which is going to
be inserted. In order to insert correctly, the function firstly checks if the tree is a \inlinecode{Leaf}, in that case the result is the \inlinecode{Node}
with the key and value given to the function. If the tree is a red node it will check if the key \inlinecode{k} is lesser or greater to insert on the
left or right subtree accordingly, if the keys are the same then the value will be updated to the most recent one, so the value \inlinecode{v}. For the
option of the tree being a black tree than it will also check if the key \inlinecode{k} is lesser or greater and will insert to the left or right
subtrees and calling the balancing function, if the keys are the same the value is updated.

This function will require the tree to satisfy the condition of \inlinecode{bst} and that exists a value in which the tree satisfies the condition
of \inlinecode{rbtree}. As a recursive function we need to say what is the value that is variating which is the tree. To ensure the correct insertion
the resulting tree needs to satisfy the condition of \inlinecode{bst}, for all values that satisfy the condition of \inlinecode{rbtree} for the original
tree the resulting tree needs to satisfy the condition of \inlinecode{almost\_rbtree} for the same value. Because of the different way we insert depending
on if the original tree was red or black, if the tree was black than the resulting tree will already satisfy the condition of \inlinecode{rbtree}. The key and
value of the new inserted node must also be present in the resulting tree and for every key \inlinecode{k'} and value \inlinecode{v\'}
that satisfies \inlinecode{memt} for the resulting tree than it is equivalent to satisfy \inlinecode{memt} for the left and right subtrees if the key \inlinecode{k\'}
is not the same as the key \inlinecode{k}, otherwise the values must be the same.


\begin{whylang}
    let rec insert (t : tree) (k : key) (v : value) : tree
        requires { bst t /\ exists n: int. rbtree n t }
        variant { t }
        ensures { bst result /\
        (forall n : int. rbtree n t -> almost_rbtree n result /\
            (is_not_red t -> rbtree n result)) /\
        memt result k v /\
        forall k':key, v':value.
        memt result k' v' <-> if k' = k then v' = v else memt t k' v' }
    = match t with
        | Leaf ->
            Node Red Leaf k v Leaf
        | Node Red l k' v' r ->
            if k < k' then Node Red (insert l k v) k' v' r
            else if k' < k then Node Red l k' v' (insert r k v)
            else (* k = k' *) Node Red l k' v r
        | Node Black l k' v' r ->
            if k < k' then lbalance (insert l k v) k' v' r
            else if k' < k then rbalance l k' v' (insert r k v)
            else (* k = k' *) Node Black l k' v r
        end
\end{whylang}

The last function of this data structure is the \inlinecode{add} that receives the tree \inlinecode{t}, the key \inlinecode{k} and
the value \inlinecode{v}. The process passes from matching the resulting tree of the function \inlinecode{insert t k v} and matching with
any type of tree and changing the color of that tree into black, if the resulting tree is a \inlinecode{Leaf} then it is an absurd because 
it is impossible because at least one node must have.

The requirements for this function is the tree \inlinecode{t} to satisfy the condition \inlinecode{bst} and that exists of value that satifies the condition
of \inlinecode{rbtree} for the \inlinecode{t}. The function must ensure that the resulting tree satisfies the condition \inlinecode{bst}, there exists a value
that the resulting tree satisfies the condition of \inlinecode{rbtree}, the key \inlinecode{k} and the value \inlinecode{v} exist inside the resulting
tree and for every key \inlinecode{k'} and value \inlinecode{v\'} that satisfies \inlinecode{memt} for the resulting tree than it is equivalent to satisfy 
\inlinecode{memt} for the left and right subtrees if the key \inlinecode{k\'} is not the same as the key \inlinecode{k}, otherwise the values must be the same.

\begin{whylang}
    let add (t : tree) (k : key) (v : value) : tree
        requires { bst t /\ exists n:int. rbtree n t }
        ensures { bst result /\ (exists n:int. rbtree n result) /\
        memt result k v /\
        forall k':key, v':value.
        memt result k' v' <-> if k' = k then v' = v else memt t k' v' }
    = match insert t k v with
        | Node _ l k' v' r -> Node Black l k' v' r
        | Leaf -> absurd
        end
\end{whylang}

The complete code is in the appendix ~\ref{subsec:Red_Black_Trees_Whyml_Og}.

\subsubsection{Extraction results}

The verified code used in this case study is taken directly from the \whythree library of examples~\ref{subsec:Red_Black_Trees_Whyml_Og}, which 
provides a fully verified \whyml implementation of the Red-Black tree data structure, including all the necessary invariants and proof obligations. 
Since this implementation originates from the \whythree library, its correctness has already been established and all proof obligations are discharged, 
satisfying the prerequisite for extraction. The code was then subjected to the extraction mechanism described in Chapter~\ref{cha:Extraction}, 
targeting \cml as the output language.

The extraction process completed fully automatically, requiring no manual corrections to either the \whyml source or the generated \cml 
output~\ref{subsec:Red_Black_Trees_Cakeml}. This result demonstrates that the corrected extraction driver is capable of handling a non-trivial 
verified data structure and producing compilable \cml code without user intervention. This is a significant indicator of the robustness of the 
corrected pipeline for programs that operate within the defined \whyml subset.

\subsubsection{Translation results}

For the translation approach, the \cml code generated by the extraction process was annotated with \gospel specifications~\ref{subsec:Red_Black_Trees_Cakeml_Gospel}, 
written manually to capture the intended correctness properties of the Red-Black tree operations. This annotated \cml code was then passed through 
the translation tool, which produced the corresponding \whyml code~\ref{subsec:Red_Black_Trees_Whyml}. 

The generated \whyml code was subsequently submitted to \whythree, where all proof obligations were discharged automatically by the available provers, 
confirming that the translation process correctly preserved the verification intent of the manually written \gospel specifications.

\section{Results Comparison}

With the corrected \whythree extraction tool and the new translation tool, while always staying in line with the
\whyml subset, it is evident from the results that both approaches work accordingly with what was expected in the
expected pipeline. Both of the codes when passing through the verification tool proceeded correctly and with no errors.

The values presented in the following tables (CITE TABLES HERE) show very similar values, making no significant changes in terms of
prover efficiency. As stated in the \ref{sec:Verification_Guarantees_and_Trust_Boundaries} there seems to be only the difference in where the 
failure points is and which code it starts with, so in a way the decision to choose one approach over another is only related in what language
the tool is starting with. If the language the code is written is \ocaml or \whyml, extraction is the best choice, if the code is written in
\cml it is better to add the \gospel specification and using the translation tool.

Regarding the \gospel specification language there is sufficient proof in the new tool created that \gospel is quite easy to integrate correctly
to a language from the ML family.